% ===================== sec_setup.tex =====================
% NOTE (revision 2): rewritten against stage1/stage3 scripts, pillar_b_*.py, stage2_pillar_d_calibration.py, stage3_pd_*.py and claim_audit_report.md. Added: model citations, decoding settings of the core run, the measurement model, the mitigation and calibration experimental designs, and the statistical procedures as run in the audit.
\section{Experimental Setup}\label{sec:setup}

\subsection{Models}\label{sec:models}
We evaluated \Nmodels{} vision-language models (Table~\ref{tab:models}). Nine are general-purpose models: Qwen2.5-VL-7B~\citep{bai2025qwen25vl}, Gemma3-12B~\citep{gemma2025gemma3}, InternVL3-8B~\citep{zhu2025internvl3}, Qwen3.5-9B~\citep{qwen2026qwen35}, Qwen3-VL-8B~\citep{bai2025qwen3vl}, MiniCPM-V-2.6~\citep{yao2024minicpmv}, Idefics3-8B-Llama3~\citep{laurencon2024idefics3}, Qwen2-VL-2B~\citep{wang2024qwen2vl} and LLaVA-1.6-Mistral-7B~\citep{liu2024llavanext}. Three were developed for medical use: LLaVA-Med v1.5 (7B)~\citep{li2023llavamed}, CheXagent (8B)~\citep{chen2024chexagent} and HuatuoGPT-Vision (7B)~\citep{chen2024huatuogptvision}. We call these three the domain-adapted models, and they form the comparison group in the specialisation ratio index. Proprietary models were not evaluated, because their use at the scale of \Nresponses{} responses was not possible within the available resources. The general-purpose models are not matched to the domain-adapted models in size, architecture or date of release, and the design contains no pair of a base model and its adapted version.
% NOTE: the notebook lists 15 planned models; three were dropped (Aya-Vision-8B, LLaVA-Interleave-Qwen-7B, Gemma4:e4b) and the reason is not recorded (reconcile_questions.md, B1). State the reason here or in the limitations.

\begin{table}[t]
\centering
\small
\begin{tabular}{llrl}
\toprule
Model & Type & Params (B) & Checkpoint \\
\midrule
Qwen2.5-VL-7B & General & 7 & \texttt{Qwen/Qwen2.5-VL-7B-Instruct} \\
Gemma3-12B & General & 12 & \texttt{google/gemma-3-12b-it} \\
InternVL3-8B & General & 8 & \texttt{OpenGVLab/InternVL3-8B} \\
Qwen3.5-9B (Omni) & General & 9 & \texttt{Qwen/Qwen3.5-9B} \\
Qwen3-VL-8B & General & 8 & \texttt{Qwen/Qwen3-VL-8B-Instruct} \\
MiniCPM-V-2.6 & General & 8.1 & \texttt{openbmb/MiniCPM-V-2\_6} \\
Idefics3-8B-Llama3 & General & 8 & \texttt{HuggingFaceM4/Idefics3-8B-Llama3} \\
Qwen2-VL-2B & General & 2 & \texttt{Qwen/Qwen2-VL-2B-Instruct} \\
LLaVA-1.6-Mistral-7B & General & 7 & \texttt{llava-hf/llava-v1.6-mistral-7b-hf} \\
HuatuoGPT-Vision-7B & Medical specialist & 7 & \texttt{FreedomIntelligence/HuatuoGPT-Vision-7B} \\
CheXagent-8B & Medical specialist & 8 & \texttt{StanfordAIMI/CheXagent-8b} \\
LLaVA-Med-v1.5-7B & Medical specialist & 7 & \texttt{microsoft/llava-med-v1.5-mistral-7b} \\
\bottomrule
\end{tabular}
\caption{The 12 evaluated VLMs, in the fixed presentation order (general models, then medical specialists).}
\label{tab:models}
\end{table}


\subsection{Inference}\label{sec:inference}
All models were run from their public checkpoints on a single GPU with 16 GB of memory, with 4-bit NF4 quantisation~\citep{dettmers2023qlora}. Quantisation was needed to fit the larger models in memory. It can change model behaviour, so absolute hallucination rates may differ from those of full-precision inference. One response was generated for each image--query pair. In the core run, responses were generated by sampling at a temperature of 0.1 with a nucleus threshold of 0.9 for most models, and not by greedy decoding. Generation stopped at the token limits given in Section~\ref{sec:queries}, so a response that reached its limit ends before the model finished.
% NOTE: confirm top-p and the decoding parameters per model against the inference notebook and the logs (the audit states top-p 0.9 for most models only).
For LLaVA-Med, a decoding defect caused a substantial share of responses to begin in mid-sentence, and Section~\ref{subsec:res_sri} reports how this affects the comparison.
% NOTE: the audit gives 32 per cent of responses; figure_manifest.md gives 50 per cent of Q1 responses. State the basis of the percentage once it is settled.

\subsection{Measurement model}\label{sec:measurement}
Claim extraction and judging were carried out by Qwen2.5-7B-Instruct, which was used with 4-bit quantisation, greedy decoding and no access to the image. The same model performs both tasks. Four of the twelve evaluated models belong to the Qwen family, and a judge may favour text from models similar to itself, so that the possibility of a bias in this direction cannot be excluded. No clinician annotation of the judge's labels is available. Section~\ref{sec:validation} reports an audit of the labels against annotations from six simulated annotators, and Section~\ref{subsec:res_robust} examines the sensitivity of the main conclusions to label noise. The first layer of the pipeline is independent of the judge and supplies the omission counts, which account for between 25 and 42 per cent of the weighted error.

\subsection{Mitigation experiments}\label{sec:setup_mitigation}
Two inference-time strategies were compared on a subset of 500 images and five models: Qwen2.5-VL and LLaVA-1.6 as general-purpose models, and CheXagent, HuatuoGPT-Vision and LLaVA-Med as domain-adapted models. The subset is a proportional stratified sample of the benchmark. The allocation to the thirteen strata uses the largest-remainder method, and images were drawn without replacement with seed 42. It gives 1{,}679 queries per model: three queries for each of the 500 images and the fourth query for the 179 normal images. Both strategies use greedy decoding and the token limits of the core run.

\paragraph{Retrieval-augmented prompting.} For each modality we wrote a set of reference snippets, which contain definitions of the findings, diagnostic categories and thresholds that appear in the reference labels, for example the dermoscopic features of each lesion category and the cup-to-disc ratios that are usual in normal and glaucomatous eyes. The snippets contain no labels of individual images and no instructions to the model. They are embedded with the all-MiniLM-L6-v2 sentence encoder~\citep{reimers2019sbert}, and the three snippets with the highest cosine similarity to the query text are placed before the original prompt under the heading ``Reference context''. The image is not available to the text encoder, so the retrieved snippets depend only on the modality, the query type and the stratum. A retrieval with a top similarity below 0.20 falls back to the original prompt. The knowledge base was authored by the authors and has not been reviewed by a clinician. The design follows retrieval-augmented generation~\citep{lewis2020rag}.

\paragraph{Contrastive decoding.} We used visual contrastive decoding~\citep{leng2024vcd}. At each decoding step the model is run on the original image and on a distorted copy to which Gaussian noise has been added following the forward process of a diffusion model, at step 500 of 1{,}000. The next-token logits are formed as $(1+\alpha)\,\ell_{\mathrm{orig}} - \alpha\,\ell_{\mathrm{dist}}$ with $\alpha=1$, restricted to the tokens whose probability under the original image is at least $0.1$ times the maximum, and the token with the highest score is taken. The noise is fixed for each image.
% NOTE: confirm alpha = 1 and beta = 0.1 as run for all five models (pillar_b_vcd.py defaults).

\paragraph{Baselines.} The comparison uses two baselines. The benchmark baseline is the response from the core run, generated by sampling at a temperature of 0.1. Because a change from sampling to greedy decoding can alter the \SMHR{} by itself, each strategy is also compared with a matched greedy baseline, taken from the greedy configuration of the calibration experiment described below. The matched baseline exists only for responses that stated a confidence, because only these were labelled in that experiment. For LLaVA-Med the corrected greedy rerun was used, and for LLaVA-1.6 the greedy responses regenerated without the token cap that was applied to part of the calibration runs.
% NOTE: T4, Fig 9 and their captions use the benchmark baseline but call it a matched greedy baseline. Fix the captions or regenerate against the matched baseline. The planned third strategy, fine-tuning, was not carried out.

\subsection{Calibration experiment}\label{sec:setup_calibration}
The same five models and the same 500-image subset were run under six decoding configurations: greedy decoding and nucleus sampling with a threshold of 0.9, each at temperatures 0, 0.3 and 0.7. A temperature has no effect under greedy decoding, so the three greedy configurations produce identical responses and four configurations are distinct. The stated confidence was read from the response, and responses without a stated confidence were excluded from the calibration analysis, which gives a coverage between 42 and 90 per cent by model and configuration. A response is labelled correct if it contains no hallucination of any of the seven classes, counting both the judged claims and the omissions of Layer 1 (the strict label). A variant that counts the judged claims only was used as a sensitivity check. Calibration error is the expected calibration error with ten equal-width bins, and the adaptive calibration error uses ten bins of equal mass. New responses were labelled by the same pipeline as the core run, and responses identical to a core response reuse its labels.

\subsection{Statistical procedures}\label{sec:stats}
Responses to different queries about the same image are not independent, so uncertainty is estimated by resampling images and not claims or responses. Intervals are 95 per cent percentile intervals with a fixed seed (42). Within a resample, all models and conditions use the same images, so that differences between them are paired. The main tables and figures (Table~\ref{tab:main} and Figures~\ref{fig:ranking}, \ref{fig:sri} and \ref{fig:mitigation}) use 2{,}000 resamples stratified by design stratum. The tests and sensitivity analyses reported in Section~\ref{sec:results} use 5{,}000 resamples of images within modality (2{,}000 for the calibration experiment). For each model we also report the 95 per cent range of its rank across resamples. For the specialisation ratio index the reference model is held fixed within each resample. Paired differences are tested with a sign-flip permutation test on the image-level contributions, with 10{,}000 to 20{,}000 permutations. Rank correlations are Spearman correlations with a Monte Carlo permutation $p$-value from 100{,}000 permutations and a Fisher interval, which for twelve models is wide. The $p$-values are adjusted by the Holm method within each of the following families: the ten pairwise comparisons among the five leading models, the twelve cells of the specialisation index, the ten comparisons of strategy and model in the mitigation experiment for each baseline, the three correlations between the metrics, and the six correlations between the modalities. Where the intervals of two models overlap, we do not claim an order between them.
% NOTE: the main tables (2,000 stratified resamples) and the audit statistics (5,000 resamples, not stratified) use different bootstrap configurations, so intervals can differ in the third decimal. Regenerate everything from one configuration before submission. Fig 6 labels t-test Holm values (0.008, 0.074, 0.017); the audit uses permutation values (0.013, 0.088, 0.024).
